%!TEX TS-program = pdflatex
%!TEX encoding = UTF-8 Unicode
% ─────────────────────────────────────────────────────────────
%  CV estilo Harvard — Alexander Daniel Rios
% ─────────────────────────────────────────────────────────────
\documentclass[9pt]{extarticle}

% ── Paquetes ─────────────────────────────────────────────────
\usepackage[T1]{fontenc}
\usepackage[utf8]{inputenc}
\usepackage[margin=0.75in, top=0.6in, bottom=0.6in]{geometry}
\usepackage{enumitem}
\usepackage{hyperref}
\usepackage{xcolor}
\usepackage{array}
\usepackage{parskip}

% ── Colores ──────────────────────────────────────────────────
\definecolor{linkcolor}{HTML}{003366}

% ── Hyperref ─────────────────────────────────────────────────
\hypersetup{colorlinks=true, urlcolor=linkcolor, linkcolor=linkcolor,
  pdfauthor={Alexander Daniel Rios}, pdftitle={CV — Alexander Daniel Rios}}

% ── Espaciado base ───────────────────────────────────────────
\pagestyle{empty}
\setlength{\parindent}{0pt}
\setlength{\parskip}{0pt}

% ── Encabezado de sección ────────────────────────────────────
\newcommand{\cvsection}[1]{%
  \vspace{6pt}%
  {\small\bfseries\uppercase{#1}}%
  \vspace{2pt}\hrule height 0.5pt\vspace{5pt}%
}

% ── Entrada de proyecto: {Título}{Fecha}{bullets}{tech}{links}
\newcommand{\project}[5]{%
  \textbf{#1} \hfill \textit{\small #2}\\[-14pt]
  \begin{itemize}[leftmargin=1.2em, itemsep=0pt, topsep=2pt, parsep=0pt]
    #3
  \end{itemize}
  \vspace{1pt}
  {\small\textbf{Tecnologías:} #4 \quad \textbf{Enlaces:} #5}
  \vspace{5pt}%
}

\newcommand{\lastproject}[5]{%
  \textbf{#1} \hfill \textit{\small #2}\\[-14pt]
  \begin{itemize}[leftmargin=1.2em, itemsep=0pt, topsep=2pt, parsep=0pt]
    #3
  \end{itemize}
  \vspace{1pt}
  {\small\textbf{Tecnologías:} #4 \quad \textbf{Enlaces:} #5}%
}

% ── Entrada de educación: {Inst}{Título}{Loc}{Fecha}{desc}
\newcommand{\eduentry}[5]{%
  \textbf{#2}, \textit{#1} \hfill \textit{\small #4}\\
  \ifx&#3&\else{\small #3}\\[-8pt]\fi
}

\newcommand{\lasteduentry}[4]{%
  \textbf{#2}, \textit{#1} \hfill \textit{\small #4}\\
  \ifx&#3&\else{\small #3}\\[1pt]\fi
  \vspace{-14pt}
}

\newcommand{\addentry}[5]{%
  \textbf{#2}, \textit{#1} \hfill \textit{\small #4}\\
  \ifx&#3&\else{\small #3}\\[1pt]\fi
  \ifx&#5&\else{\small #5\\[-8pt]}\fi
}

\newcommand{\lastaddentry}[5]{%
  \textbf{#2}, \textit{#1} \hfill \textit{\small #4}\\
  \ifx&#3&\else{\small #3}\\[1pt]\fi
  \ifx&#5&\else{\small #5\\[-14pt]}\fi
}

% ── Entrada de publicación: {Título}{Medio}{Año} ─────────────
\newcommand{\pubentry}[3]{%
  \textbf{#1} — {\small\textit{#2}} \hfill \textit{\small #3}\\[-10pt]
}

\newcommand{\lastpubentry}[3]{%
  \textbf{#1} — {\small\textit{#2}} \hfill \textit{\small #3}\\[-14pt]
}

% ─────────────────────────────────────────────────────────────
\begin{document}

% ── ENCABEZADO ───────────────────────────────────────────────
\begin{center}
  {\LARGE\bfseries Alexander Daniel Rios}\\[4pt]
  {\small
    Científico de Datos \(\cdot\) Ingeniero de Machine Learning\\[3pt]
    Buenos Aires, Argentina
    \(\mid\) \href{mailto:alexanderdaniel_rios@hotmail.com}{alexanderdaniel\_rios@hotmail.com}
    \(\mid\) +54 9 11 5773-4409\\[2pt]
    \href{https://github.com/aletbm}{github.com/aletbm}
    \(\mid\) \href{https://www.linkedin.com/in/alexander-daniel-rios}{linkedin.com/in/alexander-daniel-rios}
    \(\mid\) \href{https://www.linktr.ee/aletbm}{linktr.ee/aletbm}
  }
\end{center}

\vspace{2pt}

% ── PERFIL PROFESIONAL ───────────────────────────────────────
\cvsection{Perfil Profesional}

Científico de Datos e Ingeniero de ML con más de 4 años de experiencia en proyectos de alto
impacto. Especializado en el ciclo de vida completo de modelos de ML, desde la concepción
hasta producción, utilizando arquitecturas MLOps (MLflow, GCP, Docker). Experto en la
aplicación de Deep Learning para PLN y Visión por Computadora para resolver problemas complejos
del mundo real.

% ── PROYECTOS RELEVANTES ─────────────────────────────────────
\cvsection{Proyectos Relevantes}

\project
  {Asistente de Investigación RAG – IA para Literatura Científica}
  {Sep 2025 – Oct 2025}
  {
    \item Construí un sistema RAG sobre \textbf{más de 120k documentos}, optimizando la recuperación en \textbf{Qdrant} con \textbf{Hit Rate 0,91} y \textbf{MRR 0,78}.
    \item Automaticé el pipeline de ingestión y embeddings con \textbf{FastEmbed y Prefect}, garantizando reproducibilidad del corpus.
    \item Integré generación con \textbf{Gemini 2.5} y monitoreo en tiempo real vía \textbf{BigQuery y Grafana}.
  }
  {Python, Qdrant, Gemini, Prefect, BigQuery, Grafana, Streamlit, Docker, FastAPI}
  {\href{https://github.com/aletbm/RAG_Research_Assistant}{Repositorio}}

\project
  {Pipeline MLOps – Predicción de Enfermedades Cardiovasculares}
  {Jul 2025 – Ago 2025}
  {
    \item Construí un pipeline E2E de predicción de riesgo clínico logrando \textbf{accuracy 0,89} con un modelo KNN optimizado.
    \item Desplegué automáticamente en \textbf{GCP Cloud Run} vía Docker y Terraform, con alertas de \textit{drift} por Slack.
    \item Orquesté entrenamiento y monitoreo con \textbf{MLflow y Prefect} para la gobernanza del modelo en producción.
  }
  {MLflow, Prefect, FastAPI, Docker, Terraform, GCP, Evidently AI}
  {\href{https://github.com/aletbm/Cardiovascular-Diseases_E2E_MLPipeline}{Repositorio} \(\mid\) \href{https://cardiovascular-diseases-api-761922006747.us-east1.run.app}{API}}

\project
  {Visión por Computadora – Diagnóstico de Cáncer en Células Sanguíneas}
  {Dic 2024 – Ene 2025}
  {
    \item Diseñé una arquitectura \textbf{UNET} para segmentación y clasificación de leucemia: \textbf{ROC AUC 0,99}, \textbf{BinaryIoU 0,94}.
    \item Desarrollé una herramienta de inferencia en tiempo real para hematología, reduciendo el tiempo de análisis manual de muestras microscópicas.
    \item Desplegué con \textbf{Kubernetes} y Flask para acceso web interactivo de alta disponibilidad.
  }
  {TensorFlow, Keras, Streamlit, Flask, Docker, Kubernetes}
  {\href{https://github.com/aletbm/Blood_Cell_Cancer_Prediction}{Repositorio} \(\mid\) \href{https://bloodcellcancerprediction.streamlit.app}{App}}

\lastproject
  {PLN – Clasificación de Tweets sobre Desastres}
  {Jul 2024 – Ago 2024}
  {
    \item Construí un modelo híbrido \textbf{BERT + LSTM} para detección de crisis en tiempo real, superando las líneas base en un \textbf{12\%}.
    \item Procesé y vectoricé \textbf{más de 10k tweets} con spaCy, optimizando la limpieza de texto para entornos de datos ruidosos.
    \item Implementé un dashboard en Streamlit para visualización masiva y clasificación de alertas críticas.
  }
  {Transformers (BERT), TensorFlow, spaCy, scikit-learn, Streamlit}
  {\href{https://github.com/aletbm/NLP_with_Disaster_Tweets}{Repositorio} \(\mid\) \href{https://disastertweetsclassification.streamlit.app}{App}}

% ── EDUCACIÓN ────────────────────────────────────────────────
\cvsection{Educación}

\eduentry
  {UTN — Universidad Tecnológica Nacional}
  {Ingeniería Electrónica}
  {Buenos Aires, Argentina}
  {Abr 2014 – Actualidad (Egreso previsto 2026)}
  {}

\lasteduentry
  {UTN — Universidad Tecnológica Nacional}
  {Técnico Universitario en Electrónica}
  {Buenos Aires, Argentina}
  {Abr 2014 – Mar 2018}

% ── FORMACIÓN COMPLEMENTARIA ─────────────────────────────────
\cvsection{Formación Complementaria}

\addentry
  {DataTalks.Club}
  {\href{https://drive.google.com/file/d/1T1NxRtjeb5xgFjWrEk9oTdtdPcqPQH2f/view?usp=sharing}{Machine Learning Zoomcamp}}
  {}
  {Sep 2024 – Ene 2025}
  {Ingeniería de ML end-to-end: modelado con Scikit-Learn/TensorFlow, despliegue con Docker, Kubernetes y serverless.}

\addentry
  {DataTalks.Club}
  {\href{https://drive.google.com/file/d/1f6o5BkSsiKrAkRrmSggAwh3WxUNkDD7v/view?usp=sharing}{Machine Learning Operations Zoomcamp}}
  {}
  {May 2025 – Ago 2025}
  {MLOps: seguimiento de experimentos (MLflow), orquestación de pipelines, monitoreo (Evidently AI), IaC en GCP con Terraform.}

\lastaddentry
  {Hugging Face}
  {\href{https://drive.google.com/file/d/1Dl4U_Ee6CsA35Ysqwebf-Jcxu8C6e_KK/view?usp=sharing}{Agentes de IA}}
  {}
  {May 2025 – Jun 2025}
  {Diseño e implementación de agentes de IA con LangGraph y LlamaIndex.}

% ── ARTÍCULOS Y PUBLICACIONES ────────────────────────────────
\cvsection{Artículos y Publicaciones}

\pubentry
  {\href{https://www.linkedin.com/pulse/computer-vision-para-el-diagn\%C3\%B3stico-de-tumores-cerebrales-rios}{Visión por Computadora para el Diagnóstico de Tumores Cerebrales}}
  {LinkedIn}
  {2023}

\pubentry
  {\href{https://datatalks.club/blog/how-to-build-blood-cell-classifier-for-cancer-prediction-case-study-from-ml-zoomcamp.html}{Cómo Construir un Clasificador de Células Sanguíneas para Predicción de Cáncer}}
  {Blog de DataTalks.Club}
  {2025}

\lastpubentry
  {\href{https://volcano-camp-325.notion.site/Natural-Language-Processing-using-spaCy-TensorFlow-and-BERT-model-architecture-2455067176b38018b7defdca1005c07a}{Procesamiento de Lenguaje Natural con spaCy, TensorFlow y Arquitectura BERT}}
  {Notion}
  {2025}

% ── IDIOMAS ──────────────────────────────────────────────────
\cvsection{Idiomas}

\begin{tabular}{@{}ll}
  \textbf{Español} & Nativo \\[2pt]
  \textbf{Inglés}  & B1 Intermedio — \href{https://links.t-educationfirst.mkt4686.com/servlet/MailView?ms=NTY0Mzg4NzcS1&r=LTgzODk5MTE0NzkS1&j=MjI1NjA1NTI4NQS2&mt=1&rt=0}{Certificado EF SET} \\
\end{tabular}

\vspace{8pt}
\begin{center}
  \small\textcolor{gray}{Última actualización: \today}
\end{center}

\end{document}